\chapter{Estructura y procesamiento de datos}

Un masterfile es un archivo de texto que contiene todos los registros DNS (Resource Records) que conforman la zona que el servidor va a manejar, pero estos datos están en crudo
y necesitan de un procesamiento refinado para poder convertirse en estructuras de datos eficientes que permitan responder consultas DNS de manera rápida.
En esta seccion se explica cómo se lleva a cabo este procesamiento y las estructuras de datos utilizadas para almacenar la información de la zona DNS en memoria.

\section{Catalogo}
Estructura que contiene un \texttt{Hashmap} donde cada llave es un string con el nombre de la zona y el valor es un puntero a la estructura con los datos reales.
Es basicamente un índice que permite acceder rápidamente a los datos de cada zona, y es especialmente útil cuando el servidor maneja múltiples zonas DNS. También
es útil para saber si una consulta pertenece a una zona manejada por el servidor o no, ya que si el dominio consultado no es parte de ninguna
zona del catalogo, entonces se puede responder inmediatamente con un \texttt{REFUSED}.

El valor de retorno en el \texttt{Hashmap} está envuelto en un \verb|Arc<RwLock<>>| para permitir el acceso concurrente a los datos de la zona desde múltiples threads.

%salto de linea para que se vea mejor
\begin{lstlisting}
    pub struct Catalog {
        zones: HashMap<String, Arc<RwLock<DataStructure>>>, 
    }
\end{lstlisting}

\section{Hash anidado de Resource Records}

